\documentclass[11pt,a4paper]{article}
\usepackage[spanish,es-noquoting]{babel}
\usepackage[top=2.5cm, bottom=2.5cm, left=2.5cm, right=2.5cm]{geometry}
\usepackage{fontspec}
\usepackage{amsmath}
\usepackage[version=4]{mhchem}
\usepackage{enumitem}
\usepackage{microtype}

\setmainfont{Calibri}
\linespread{1.2}

\begin{document}
\pagestyle{empty}

% PÁGINA 1: EJERCICIO 1
\begin{center}
    {\LARGE \textbf{Química: Estequiometría y Rendimiento de Reacciones}}\\[4pt]
    {\large \textbf{Ficha de Trabajo | Ejercicio 1: Combustión de Hidrocarburos}}
\end{center}

\vspace{0.8cm}
\noindent\textbf{1.} En una caldera de calefacción se queman $88\text{ g}$ de gas propano (\ce{C3H8}) con una cantidad suficiente de gas oxígeno (\ce{O2}) del aire, produciéndose dióxido de carbono (\ce{CO2}) y agua en estado líquido (\ce{H2O}) según la siguiente reacción sin ajustar:
\[
\ce{C3H8 (g) + O2 (g) -> CO2 (g) + H2O (l)}
\]

\begin{enumerate}[label=\textbf{\alph*)}, leftmargin=2em, itemsep=12pt]
    \item Escribe y ajusta la ecuación química de la combustión.
    \item Calcula los moles y la masa teórica máxima de dióxido de carbono (\ce{CO2}) que se producirían con un rendimiento del $100\%$.
    \item Sabiendo que las condiciones del quemador hacen que la reacción tenga un rendimiento real del $85\%$, calcula la masa real de \ce{CO2} que se expulsa a la atmósfera.
    \item Determina el volumen real de gas \ce{CO2} desprendido, medido en Condiciones Normales ($\text{C.N.}$: $1\text{ atm}$, $0\,^\circ\text{C}$, donde $1\text{ mol} = 22{,}4\text{ L}$).
\end{enumerate}

\vspace{0.5cm}
\noindent\textbf{Datos:} Masas atómicas ($\text{u}$): $\mathrm{H} = 1$; $\mathrm{C} = 12$; $\mathrm{O} = 16$.

\newpage

% PÁGINA 2: EJERCICIO 2
\begin{center}
    {\LARGE \textbf{Química: Estequiometría y Rendimiento de Reacciones}}\\[4pt]
    {\large \textbf{Ficha de Trabajo | Ejercicio 2: Proceso Metalúrgico de Reducción}}
\end{center}

\vspace{0.8cm}
\noindent\textbf{2.} En un alto horno siderúrgico se tratan $320\text{ g}$ de óxido de hierro (III) (\ce{Fe2O3}) puro mediante una corriente de monóxido de carbono (\ce{CO}) gaseoso para obtener hierro metálico (\ce{Fe}) y dióxido de carbono (\ce{CO2}) según la ecuación química sin ajustar:
\[
\ce{Fe2O3 (s) + CO (g) -> Fe (s) + CO2 (g)}
\]

\begin{enumerate}[label=\textbf{\alph*)}, leftmargin=2em, itemsep=12pt]
    \item Escribe y ajusta la ecuación química del proceso de reducción.
    \item Calcula la masa teórica máxima de hierro puro (\ce{Fe}) que se debería obtener si el rendimiento fuese del $100\%$.
    \item Si al finalizar la reacción se han obtenido experimentalmente en el crisol $179{,}2\text{ g}$ de hierro metálico, calcula el porcentaje de rendimiento ($\%$) del proceso.
    \item Calcula la masa real de dióxido de carbono (\ce{CO2}) que se ha generado en el alto horno durante la operación.
\end{enumerate}

\vspace{0.5cm}
\noindent\textbf{Datos:} Masas atómicas ($\text{u}$): $\mathrm{C} = 12$; $\mathrm{O} = 16$; $\mathrm{Fe} = 56$.

\end{document}
